\begin{exercisechunk}
\item \textbf{Pooling Bradley--Terry evaluators.}
\label[exercise]{ex:l19-pooled-evaluators}
\exerciseprofile{2}{1}{1}{\excalculation\quad\exinterpretation}{%
Show that pooling correctly specified Bradley--Terry evaluators can produce
preferences that no Bradley--Terry model represents.}
Fix a prompt with responses $a,b,c$. For each comparison, draw an evaluator
type $T$ uniformly from $\{1,2,3\}$, independently of the presented pair.
The three types have score vectors
\[
 \begin{pmatrix}r_1(a)&r_1(b)&r_1(c)\\
                 r_2(a)&r_2(b)&r_2(c)\\
                 r_3(a)&r_3(b)&r_3(c)\end{pmatrix}
 =\begin{pmatrix}\log9&\log3&0\\
                 0&\log9&\log3\\
                 \log3&0&\log9\end{pmatrix}.
\]
Conditional on type $T=t$ and the ordered pair $(y,y')$, the evaluator
selects $y$ with probability $\sigma(r_t(y)-r_t(y'))$, where
$\sigma(z)=(1+e^{-z})^{-1}$. The pooled probability averages over the
unobserved evaluator type:
$p(y,y')=\tfrac13\sum_{t=1}^3\sigma(r_t(y)-r_t(y'))$.
\begin{enumerate}[(a)]
\item Show that $p(a,b)=p(b,c)=p(c,a)=8/15$.
\item Show that no scalar Bradley--Terry score function on $\{a,b,c\}$
represents these pooled probabilities. Identify the violated cycle
restriction and explain why correct specification for each evaluator
type does not imply correct specification after the type is omitted.
\end{enumerate}
\end{exercisechunk}
